\documentclass[11pt]{article}
\usepackage[T1]{fontenc}
\usepackage{lmodern,microtype,amsmath,amssymb,amsthm,booktabs}
\usepackage[a4paper,top=24mm,bottom=24mm,left=27mm,right=27mm]{geometry}
\usepackage{xcolor,fancyhdr,enumitem}
\definecolor{ink}{HTML}{193F51}
\usepackage[colorlinks=true,linkcolor=ink,citecolor=ink,urlcolor=ink,pdftitle={Binary Machin-type formulae and the domain of Davis's 2-adic binomial sum}]{hyperref}
\newtheorem{theorem}{Theorem}[section]
\newtheorem{lemma}[theorem]{Lemma}
\newtheorem{proposition}[theorem]{Proposition}
\newtheorem{corollary}[theorem]{Corollary}
\setcounter{tocdepth}{1}
\setlength{\emergencystretch}{2em}
\pagestyle{fancy}\fancyhf{}\lhead{\small\textsc{Arctangent relations and reciprocal binomial sums}}\rhead{\small Research manuscript}\cfoot{\thepage}\setlength{\headheight}{14pt}
\begin{document}
\begin{center}
{\small\color{ink}\textsc{Number theory \enspace / \enspace Research manuscript}}\\[10pt]
{\LARGE\bfseries Binary Machin-type formulae and the domain\\[3pt] of Davis's $2$-adic binomial sum}\\[8pt]
{\small November 2026}
\end{center}
\begin{abstract}
We answer two questions of Borwein, Borwein and Galway about binary Machin-type formulae. The zero relations among their minimal arctangent generators form a two-dimensional space. Membership in their logarithmic family is decidable for every positive algebraic argument, with the Aurifeuillian generators included. The proof uses separate primitive divisors of the two factors of $2^{2k}+1$ at each large odd level. We also determine the domain of Davis's sum of reciprocal binomial coefficients along canonical dyadic truncations. If $-1-x$ has infinitely many set bits at positions $d_i$ and $R_i$ is its prefix through $d_i$, convergence is equivalent to $R_i-d_{i+1}\to+\infty$. Isolating one term in the difference between successive truncations proves necessity. We establish both of Davis's finite valuation conjectures with all equality cases. At every nonintegral rational argument in $\mathbb Z_2$ the limit is transcendental. Its linear relations at rational arguments are determined by their classes modulo $\mathbb Z$.
\end{abstract}
\tableofcontents
\section*{The original questions}
In 2004, Borwein, Borwein and Galway gave two identities for binary digit extraction and asked whether any further independent relations exist. They also asked for an exclusion test allowing both of their logarithmic generator families \cite[Appendix D, Problems 2 and 8]{BBG}. In 2012 Davis found both convergent and divergent classes of inputs for his reciprocal-binomial sum \cite{Davis12}. He then formulated two uniform valuation conjectures \cite{Davis14}.

Item 022 of the corpus \cite{Corpus} uses a prime of prescribed order to isolate one term. Here we ask whether two primes can distinguish two generators with the same highest index. Items 139 and 169 suggest examining binary carries and the size of gaps between successive digits. Items 176 and 180 suggest the Fourier calculation of the conjugate logarithms in Section~\ref{tr:section}. The two classifications have independent proofs.
\clearpage
\section{Zero relations among binary Machin-type arctangents}
\label{bbp:section}
Borwein, Borwein and Galway asked for the number of independent
binary Machin-type BBP arctangent zero relations in Appendix~D,
Problem~2 of their 2004 paper \cite{BBG}. Their Definition~1 uses the real
principal arguments
\[
 \alpha_m=\arg(2^m-1-i),\qquad
 \beta_m=\arg(2^m-i) \qquad(m\ge1).
\]
The formal coefficients of these generators have finite support and
belong to $\mathbb Q$. A zero relation has real value exactly zero.
The authors give two independent relations in their equations
(35)--(36).

\begin{theorem}\label{bbp:main}
The space of binary Machin-type BBP arctangent zero relations has
dimension two. A basis consists of the original relations
\[
    \alpha_1+\alpha_3-2\beta_1=0,\qquad
    \alpha_1-2\alpha_2-\alpha_3=0.
\]
\end{theorem}

\begin{proof}
For an odd prime $p\equiv1\pmod4$, choose one Gaussian prime
$\mathfrak p$ above $p$. On $\mathbb Q(i)^\times$ set
\[
      D_p(z)=v_{\mathfrak p}(z)-v_{\overline{\mathfrak p}}(z).
\]
This additive valuation difference vanishes on $\mathbb Q^\times$.
Clear the denominators of a putative rational relation among principal
arguments. Its Gaussian product has argument zero modulo $2\pi$, hence
is a positive rational number, since $\mathbb Q(i)\cap\mathbb R=\mathbb Q$.
Consequently every $D_p$ vanishes on this product.
For $z=u-i$ with integral $u$, any odd prime $p$ dividing $N(z)=u^2+1$
gives $D_p(z)\ne0$. Indeed, if both prime ideals above $p$ divided $z$,
then $p$ would divide both Gaussian integer coordinates of $z$.

Assign level $k$ to $\beta_k$ and level $2m-1$ to $\alpha_m$.
At odd level $k=2m-1$ put
\[
 A_k=2^k-2^{(k+1)/2}+1,\qquad B_k=2^k+2^{(k+1)/2}+1.
\]
Squaring the Gaussian moduli gives
\[
 N(2^k-i)=2^{2k}+1=A_kB_k,\qquad
 N(2^m-1-i)=2A_k .
\]
The factors $A_k,B_k$ are coprime because they are odd and their
difference is a power of two. For a prime $p$ with
$\operatorname{ord}_p(2)=4k$, no Gaussian generator of a smaller
level has norm divisible by $p$. In fact, the odd part of each
level-$h$ norm divides $2^{2h}+1$, which would force $4k\mid4h$.
Also $p\equiv1\pmod4$.

Delete $\alpha_2$ and $\alpha_3$ for the moment. We prove that all
remaining generators are independent. Consider the highest level $k>1$
of a nontrivial finite relation. For even $k$, the Bang--Zsigmondy theorem \cite[Theorem~1]{Grech} gives a prime of order $4k$ for $2$ since $4k$ is a multiple
of eight. It divides $2^{2k}+1$. Only $\beta_k$ is present at this
level, so its coefficient must vanish by $D_p$.
For odd $k\ge7$, the Suzuki-group primitive-divisor theorem
of Alp\"oge--Katz--Navarro--O'Brien--Tiep (Theorem~3.4),
also proved by Grechkoseeva (Theorem~2) \cite{AKNOT,Grech},
gives a prime $p_+$ of order $4k$ dividing $B_k$ and a prime
$p_-$ of order $4k$ dividing $A_k$. Applying $D_{p_+}$ first
annihilates every term except $\beta_k$, forcing its coefficient to
vanish. Applying $D_{p_-}$ next forces the coefficient of
$\alpha_{(k+1)/2}$ to vanish. At levels $k=3,5$ only $\beta_k$
remains, and the primes $13,41$ have orders $12,20$, respectively.
This rules out a nonzero coefficient of highest level $k>1$.
At level one, $N(2-i)=5$ and $N(1-i)=2$, so $D_5$ forces the
coefficient of $\beta_1$ to vanish. Then $\alpha_1=-\pi/4\ne0$
forces the remaining coefficient to vanish.

For completeness write $v=\arctan(1/2)$ and $w=\pi/4$. Then
$\beta_1=-v$, $\alpha_1=-w$, $\alpha_2=v-w$ and
$\alpha_3=w-2v$. The last two identities follow from the tangent
addition formula with $\tan(w-v)=1/3$ and $\tan(2v-w)=1/7$,
with all four angles in their principal ranges. Substitution gives
both displayed zero relations. Their coefficient matrix on
$(\alpha_2,\alpha_3)$ is invertible. Subtracting a unique rational
combination of them therefore deletes these two generators from any
relation. Independence of the remaining generators completes the proof.
\end{proof}

\paragraph{The enlarged presentation and integer relations.}
On page 904 the authors also allow the additional formal symbol
$\gamma_m=\arg(2^m+1+i)$. Their change of generators is
$\gamma_m+\alpha_m-\beta_{2m-1}=0$: multiplication of the first two Gaussian numbers gives $2(2^{2m-1}-i)$, with no branch ambiguity for the indicated angles. Each added relation contains its own $\gamma_m$ with coefficient one. Eliminating all such symbols identifies the quotient kernel with the two-dimensional kernel of Theorem~\ref{bbp:main}. If every redundant symbol is counted separately, its unquotiented kernel therefore has countably infinite dimension. For example, the two presentations of the angle in Appendix D, Problem 1 give $\gamma_1+\alpha_2=0$, which reduces to $\beta_1-\alpha_1+\alpha_2=0$. In the minimal presentation the integral relations have basis $\alpha_2-\alpha_1+\beta_1$ and $\alpha_3+\alpha_1-2\beta_1$, since these eliminate the two discarded coordinates over $\mathbb Z$.

\begin{corollary}[Deciding membership for rational slopes]\label{bbp:decide}
Let $z=X+iY$ be a nonzero Gaussian integer and put $N=X^2+Y^2$.
To decide if $\arg(z)$ admits a rational binary Machin-type BBP
formula, it suffices to use generators of level at most
\[
 H(z)=\max\left(\{5\}\cup
  \{\lfloor\operatorname{ord}_p(2)/4\rfloor:
                            p\text{ an odd prime divisor of }N\}\right).
\]
In particular $\arctan(2/7)$, $\arctan(4/9)$ and $\arctan(5/9)$
have no such formula.
\end{corollary}
\begin{proof}
For a level $k>H(z)$, the primitive primes used in
the proof of Theorem~\ref{bbp:main} do not divide $N$.
Thus comparison with $D_p(z)$ again removes all generators at the
largest such level. This proves the bound. For the finite decision,
compute the $D_p$ of the remaining Gaussian generators and of $z$,
using all split rational primes dividing any of their norms. Solve
the resulting rational linear system for a representation of the
valuation vector of $z$. If there is a rational solution, clear
denominators. Dividing the corresponding power of $z$ by the Gaussian
product gives a number $w$ with all split-prime valuation
differences zero. Then $w/\overline w$ has zero valuation at every
Gaussian prime because conjugation fixes the inert and ramified primes.
Thus $w/\overline w$ is a Gaussian unit, and
$\arg(w)\in(\pi/4)\mathbb Z$. Since $\alpha_1=-\pi/4$, the residual
angle and any multiple of $2\pi$ can be absorbed. This gives an exact
real formula. Absence of a rational solution precludes one.

For $7+2i$ and $9+5i$ the norms are $53$ and $2\cdot53$,
and $\operatorname{ord}_{53}(2)=52$. At the largest possible level
$13$ the norm of the $\beta$ generator equals
$5\cdot53\cdot157\cdot1613$, and that of the $\alpha$ generator is
$2\cdot 5\cdot1613$. The prime $157$ also has order $52$, so the
$\beta$ coefficient vanishes. No remaining generator accounts for
$D_{53}(z)\ne0$. For $9+4i$ the norm is $97$. At level $12$ the
single generator $\beta_{12}$ has norm $97\cdot257\cdot673$.
Both $97$ and $673$ have order $48$, so first $D_{673}$ removes
$\beta_{12}$ and then $D_{97}$ rules out the target. To check the orders,
$2^{26}\equiv-1$ modulo $53,157$ and $2^4=16$ in each field.
Also $2^{24}\equiv-1$ modulo $97,673$, while $2^{16}$ is $61,255$
respectively.
\end{proof}


\subsection{The Aurifeuillian exclusion problem for logarithms}
Appendix D, Problem~8, of \cite{BBG} asks for an exclusion criterion that retains the Aurifeuillian logarithms. Their Definition~3 and equation~(43) give
\[
 g_m=\log(1-2^{-m}),\qquad
 a_s=\log\lvert 1-(1+i)2^{-s}\rvert
     =\tfrac12\log\bigl(A_{2s-1}/2^{2s-1}\bigr),\qquad m,s\ge1,
\]
where $A_j=2^j-2^{(j+1)/2}+1$ for odd $j$. The ordinary generator $g_m$ has level $m$, and $a_s$ has level $4(2s-1)$. The authors used a primitive-divisor cutoff for their non-Aurifeuillian logarithms in Theorem~6. Two separate primitive divisors at a shared level give the following decision criterion.

\begin{theorem}\label{bbg:log-exclusion}
For a positive rational $x$, put
\[
 L=\max\bigl(\{6\}\cup
   \{\operatorname{ord}_p(2):p\text{ odd prime},\ v_p(x)\ne0\}\bigr).
\]
Then $\log x$ has a rational linear binary Machin-type BBP logarithm formula, with both families of Definition~3 allowed, if and only if it lies in the rational span of the generators of levels at most $L$. Membership in that finite span is decidable by rational linear algebra on ordinary prime valuations.
\end{theorem}
\begin{proof}
Write $B_j=2^j+2^{(j+1)/2}+1$ for odd $j$, so that $A_jB_j=2^{2j}+1$. For a generator of level $d$, the odd prime factors of its squared argument divide $2^d-1$. Thus a prime of order $m$ cannot occur at any smaller level. At a hypothetical largest level $m>L$, primes of order $m$ are absent from $x$ by definition. At a level with only $g_m$, Bang--Zsigmondy supplies a prime of order $m$, since $m>6$. Its valuation forces the coefficient of $g_m$ to vanish.

If $m=4j$ with $j>5$ odd, Grechkoseeva supplies primes of order $m$ in both $A_j$ and $B_j$. The factors are coprime because they are odd and their difference is a power of two. Valuation at a prime in $B_j$ first kills the coefficient of $g_m$, and valuation at a prime in $A_j$ then kills the coefficient of the other generator. These deductions apply to logarithmic relations after clearing rational denominators, exponentiating, and squaring to remove the positive square roots.

The two exceptional shared levels above six are $12$ and $20$. Their complementary factors $B_3=13$ and $B_5=41$ have orders $12$ and $20$, respectively: $2^6\equiv-1\pmod{13}$ with $2^4\not\equiv1\pmod{13}$, and $2^{10}\equiv-1\pmod{41}$ with $2^4\not\equiv1\pmod{41}$. Thus $g_m$ is eliminated exactly as before. The remaining generators $a_2,a_3$ have the smaller-level presentations
\[
 2a_2=g_4-g_2+g_1,\qquad
 2a_3=2g_4-2g_2+g_1,
\]
from $A_3=5$ and $A_5=25$. Successive descent terminates. The reverse implication follows by inclusion. Finally, each $g_m$ is the logarithm of a positive rational, and twice each $a_s$ is too. Equality of two positive rationals is equivalent to equality at every prime valuation. Factoring the finitely many numerators and denominators and testing their valuation columns therefore decides membership.
\end{proof}

For example, $\operatorname{ord}_{47}(2)=\operatorname{ord}_{178481}(2)=23$ and $2^{23}-1=47\cdot178481$, a factorization already used in Section~3 of Borwein, Borwein and Galway. Neither prime can occur in a generator below level $23$. The level-$23$ column is $g_{23}$ and contains both, so the valuation at $178481$ forces its coefficient to vanish in any putative expression for $\log47$. This excludes $\log47$ even with the Aurifeuillian generators. The authors already observed that $A_7=113$ at level $28$ gives $\log113=2a_4-7g_1$, illustrating why a second valuation at a shared level is necessary.

The criterion also decides membership for every positive real algebraic $x$.
A formula forces $x^d\in\mathbb Q_{>0}$ for some $d\ge1$ by clearing rational coefficients and doubling. Conversely a formula for $\log(x^d)$ gives one for $\log x$. To decide the rational-power condition, choose a normal closure $K$ of $\mathbb Q(x)$, of degree $D$. Some power of $x$ is rational exactly when every quotient $\sigma(x)/x$, for $\sigma\in\operatorname{Gal}(K/\mathbb Q)$, is a root of unity. In the affirmative case take $d$ to be the least common multiple of their orders. The test is finite: an order $m$ root of unity in $K$ has $\varphi(m)\le D$ and $m\le2D^2$. The second inequality follows by multiplying $\varphi(p^a)^2/p^a=p^{a-2}(p-1)^2$ over prime powers, whose only factor below one is $1/2$ for $p^a=2$.

\begin{corollary}\label{bbp:log23}
Neither $\log 23$ nor $\log 89$ has a binary Machin-type logarithm formula with the full set of generators in Definition~3 of \cite{BBG}.
\end{corollary}
\begin{proof}
Since $2^{11}-1=23\cdot89$ and $11$ is prime, both primes have order $11$. The cutoff in Theorem~\ref{bbg:log-exclusion} is $11$. Only $g_{11}$ contains either prime through this level and its valuations at the two primes agree. The target valuation vectors have unequal entries there.
\end{proof}
The original Theorem~4 of \cite{BBG} excluded ordinary generators alone. Its later remark (i) explicitly leaves the Aurifeuillian possibility for $\log 23$ open.

\begin{corollary}[Relations among the real logarithmic generators]
All rational linear relations among $g_m,a_s$ are generated by the four identities
\begin{align*}
 g_6&=2g_2+g_3-g_1,&2a_1&=g_1,\\
 2a_2&=g_4-g_2+g_1,&2a_3&=2g_4-2g_2+g_1.
\end{align*}
\end{corollary}
\begin{proof}
The identities follow from $2^6-1=3^2\cdot7$ and $A_1=1,A_3=5,A_5=25$. Eliminate $g_6,a_1,a_2,a_3$ using their four distinct coordinates. Above level six the descending argument in Theorem~\ref{bbg:log-exclusion} eliminates every remaining coefficient, including those of the sole surviving generators at levels 12 and 20. At levels five down to two the primes $31,5,7,3$ successively eliminate $g_5,g_4,g_3,g_2$. Finally $g_1=-\log 2\ne0$. 
\end{proof}

\section{An auxiliary problem on moving denominators}\label{sec:aux}
Fix a prime $p$. Write $v=v_p$ with $v(0)=+\infty$ and $|y|_p=p^{-v(y)}$.
For $x\in\mathbb Z_p$, let $x_j$ be its least residue in
$[0,p^j)$, starting with $x_0=0$. Take coefficients $a_r\in\mathbb Q_p$
such that
\begin{equation}\label{eq:weights}
                   v(a_r)\ge\eta r-B\qquad(r\ge1)
\end{equation}
for some real $\eta>0$ and $B$. For integers $n\ge0$ put
\[
 H_a(n)=-(n+1)\sum_{r\ge1}\frac{a_r}{n+1+r}.
\]
This converges since $v(n+1+r)\le\log_p(n+1+r)$. The auxiliary problem is to decide from $a$ and $x$ whether $H_a(x_j)$ is Cauchy. Substituting $x$ formally in the displayed series suggests a sufficient test. Necessity requires controlling all remaining denominators at a truncation where one of them is unusually small in the $p$-adic metric.

\begin{theorem}[Convergence of the moving-denominator sums]\label{th:aux}
Set $c=-1-x$ and omit terms for which $a_r=0$. If $c$ is a positive
integer and $a_c\ne0$, the sequence $H_a(x_j)$ diverges.
In all other cases it converges if and only if
\begin{equation}\label{eq:weightcriterion}
              v(a_r)-v(r-c)\longrightarrow+\infty
              \quad(r\longrightarrow\infty,\ a_r\ne0).
\end{equation}
The condition holds vacuously for finite support. On the domain of
convergence the limit is
\[
                 c\sum_{r\ge1,\ a_r\ne0}\frac{a_r}{r-c}.
\]
\end{theorem}
\begin{proof}
For every retained index with $r\ne c$ set $w_r=v(r-c)$. There is a uniform upper bound
\begin{equation}\label{eq:denom}
                 v(x_j+1+r)\le 1+\max\{w_r,\log_p(1+r)\}.
\end{equation}
When $w_r<j$, the left side equals $w_r$ by congruence. Otherwise $x_j+1+r\le p^j+r\le p\max\{p^{w_r},1+r\}$. Under \eqref{eq:weightcriterion}, both $v(a_r)-w_r$ and $v(a_r)-\log_p(1+r)$ tend to infinity along the support. Thus the tails are uniformly small for every $j$. Each term converges to its asserted limit and one may pass to the limit in the sum.

For necessity we compare successive truncations. For positive integers $N,t$ direct subtraction gives
\begin{equation}\label{eq:pair}
 H_a(N+t-1)-H_a(N-1)=-t\sum_{s\ge1}\frac{a_s s}{(N+s)(N+t+s)}.
\end{equation}
Suppose there is no retained pole and the test fails. Choose retained
$r\to\infty$ with $w=v(r-c)\ge v(a_r)-C$ for one fixed $C$. In particular $w\ge\eta r-O(1)$. The case $c=0$ cannot occur, by $v(r)\le\log_p r$. Thus $v(r)=v(c)$ for all sufficiently large selected $r$.
Set $j=w-1$ and $T=p^j$. Eventually $0<r<T-1$. Since $c\equiv r\pmod{pT}$ we have
\[
               x_j=T-1-r,\qquad x_{j+1}=pT-1-r .
\]
Use \eqref{eq:pair} with $N=T-r$ and $t=(p-1)T$. The $s=r$ summand has denominators $T,pT$ and $v(t)=j$, giving valuation
\begin{equation}\label{eq:badterm}
              v(a_r)+v(r)-w\le C+v(c).
\end{equation}
For $1\le s<T$ with $s\ne r$, both denominators have valuation
$v(s-r)\le\log_p\max(s,r)$. Hence each other term in this range has valuation at least
\[
                 j+\eta s-B-2\log_p\max(s,r).
\]
For $s\le r$ this is at least $j-B-2\log_p r$. For $s>r$ use
$\inf_{s\ge1}(\eta s-2\log_p s)>-\infty$. For $s\ge T$ both
denominators are positive and at most $(p+1)s$, giving the lower bound
$j+\eta s-B-2\log_p((p+1)s)$. Together these bounds tend uniformly to
infinity as the selected $r$ grows, since $j\ge\eta r-O(1)$.
The chosen summand therefore has uniquely smallest valuation. The
increments have bounded-above valuations and $H_a(x_j)$ is not Cauchy.

Finally take $c=m\in\mathbb Z_{>0}$ with $a_m\ne0$.
Use the same pair of truncations with $r=m$ fixed and $j$ tending to infinity.
The $m$th term has valuation $v(a_m)+v(m)-j-1$, and the other terms
again tend uniformly to zero. This proves the pole assertion. When $a_m=0$, all retained denominators are covered by the preceding sufficiency argument, with $v(r-m)=O(\log r)$ for $r>m$.
\end{proof}

\section{The exact domain for Davis's sum}\label{sec:domain}
From now on $p=2$ and $v=v_2$. Write $f(n)=\sum_{k=0}^n\binom nk^{-1}$ and $F(x)=\lim_{j\to\infty}f(x_j)$ when it exists.
For $N=n+1$ the classical finite identity \cite{Rockett} is
\begin{equation}\label{eq:finite}
             f(n)=\frac{N}{2^N}\sum_{k=1}^N\frac{2^k}{k}.
\end{equation}
Indeed, symmetry of $\binom nj$ gives $f(n+1)=1+(n+2)f(n)/(2(n+1))$ upon summing $\binom{n+1}{j}^{-1}=(n+1-j)\binom nj^{-1}/(n+1)$ over $0\le j\le n$ and adding the final term. The right side of \eqref{eq:finite} satisfies the same recurrence and has initial value $1$.
The $2$-adic logarithm yields $\sum_{k\ge1}2^k/k=-\log(-1)=0$.
Completing the sum as in Cohen \cite[Lemmas 2 and 4]{Cohen} gives the exact identity
\begin{equation}\label{eq:tail}
          f(n)=-N\sum_{r\ge1}\frac{2^r}{N+r}=H_{(2^r)}(n).
\end{equation}

\begin{theorem}[Domain of $2$-definability]\label{th:domain}
For $x\in\mathbb Z_2$ let $c=-1-x$. If $c\in\mathbb Z_{>0}$ then $F(x)$ does not exist. Otherwise
\begin{equation}\label{eq:iff}
 F(x)\text{ exists }\quad\Longleftrightarrow\quad
          r-v(r-c)\longrightarrow+\infty
            \qquad(r\longrightarrow\infty,\ r\in\mathbb Z_{>0}).
\end{equation}
On this domain, $F(x)=c\sum_{r\ge1}2^r/(r-c)$, in particular $F(-1)=0$.
\end{theorem}
\begin{proof}
Specialize Theorem~\ref{th:aux} to $a_r=2^r$ in \eqref{eq:tail}.
\end{proof}
Here $x\mapsto c$ is an involution of $\mathbb Z_2$ and is used only to index denominators. Identity~\eqref{eq:tail} holds at every original integer truncation before limits are taken.

\begin{corollary}[Binary form]\label{th:digits}
Suppose $c$ has infinitely many ones at positions $0\le d_0<d_1<\cdots$.
Set $R_i=\sum_{k=0}^i2^{d_k}$. Then
\[
                  F(x)\text{ exists}\quad\Longleftrightarrow\quad
                               R_i-d_{i+1}\longrightarrow+\infty .
\]
The digits of $c$ are exactly the complements of the digits of $x$.
\end{corollary}
\begin{proof}
We have $v(R_i-c)=d_{i+1}$. For any $R_i\le r<R_{i+1}$,
congruence modulo $2^{d_{i+1}+1}$ with $c$ would require
$r\ge R_{i+1}$, its least such residue. Thus $v(r-c)\le d_{i+1}$,
and the minimum of $r-v(r-c)$ on this interval is $R_i-d_{i+1}$.
\end{proof}

\paragraph{Two outcomes at the same radius.}
For $c$ with infinitely many ones let $\rho(c)$ denote the radius of convergence of $\sum_{r\ge1}z^r/(r-c)$ in $\mathbb C_2$. The coefficient formula gives
$\rho(c)=2^{-\limsup_{r\to\infty}v(r-c)/r}$. The prefix argument also gives $\limsup_r v(r-c)/r=\limsup_i d_{i+1}/R_i$. This radius is the classical type of $c$ in the terminology of \cite{Kedlaya}. Start at $d_0=3$. Recursively choose either $d_{i+1}=R_i$ for all $i$, or $d_{i+1}=R_i-i$ for all $i$. Both rules give strictly increasing positions (by induction $d_i\ge i+3$) and both yield $\rho(c)=1/2=|2|_2$. In the first example the truncations diverge. In the second they converge. The additive difference in Corollary~\ref{th:digits} decides the common boundary.

\section{Size of the domain and of its complement}\label{sec:metric}

\begin{proposition}\label{dav:measure}
If $x\notin\{-2,-3,\ldots\}$ and $w_r=v(x+1+r)=o(r)$, there is a constant $C_x$ such that
$v(F(x)-f(x_j))\ge j-C_x$ for every $j$. The hypothesis holds for every algebraic $x\in\mathbb Z_2$ outside those negative integers and for Haar-almost every $x$.
\end{proposition}
\begin{proof}
By Theorem~\ref{th:domain} the value exists. Subtract the convergent tails at $x$ and $x_j$ by the algebra of \eqref{eq:pair}. After factoring out $x-x_j$, each summand has valuation at least
\[
      r+v(r)-w_r-1-\max\{w_r,\log_2(1+r)\}.
\]
The right side tends to $+\infty$ independently of $j$, hence is bounded below by some $-C_x$. For algebraic $x$ let $P\in\mathbb Z[X]$ be its irreducible polynomial. The integer $P(-1-r)$ is nonzero and is divisible in $\mathbb Z_2$ by $x+1+r$. Thus $w_r\le\log_2|P(-1-r)|=O_P(\log(2+r))$. For the Haar assertion, the ball $w_r\ge\lceil\epsilon r\rceil$ has probability $2^{-\lceil\epsilon r\rceil}$. These probabilities sum to a finite number. Apply the first Borel--Cantelli lemma for $\epsilon=1/q$, $q=1,2,\ldots$.
\end{proof}

\begin{theorem}\label{th:metric}
The divergence set $E$ is comeager in $\mathbb Z_2$ and has Hausdorff dimension zero. In every nonempty $2$-adic ball $U$ its exact logarithmic Hausdorff threshold is
\[
  \mathcal H^{h_s}(E\cap U)=
  \begin{cases}
    \infty,&0<s\le1,\\
    0,&s>1,
  \end{cases}
  \qquad h_s(\rho)=\bigl(\log_2(1/\rho)\bigr)^{-s}
       \quad(0<\rho\le1/2).
\]
\end{theorem}
\begin{proof}
Work in the isometric coordinate $c=-1-x$. Apart from the countable set of positive integers, \eqref{eq:iff} identifies $E$ with
\[
 \bigcup_{C=0}^{\infty}\limsup_{r\to\infty}
                  \{c:v(c-r)\ge r-C\}.
\]
The limsup for $C=0$ is a dense $G_\delta$. Each tail union is open and meets every residue class because its centers include arbitrarily large integers in that class. The Baire category theorem gives density. The displayed cover also shows that every positive power of the diameter has finite sum in each tail, giving dimension zero. For $s>1$ the gauge sum is bounded by a constant times $\sum_{r>C+1}(r-C)^{-s}$, giving zero $h_s$-measure.

For the lower bound take balls centered at $r$ with radius $2^{-r}$. Their limsup lies in $E$ for either open or closed ball conventions. Under the general mass transference principle of Beresnevich and Velani \cite[Theorem 3]{BV} with $g(\rho)=\rho$ and $h=h_1$, the enlarged balls have radius $h_1(2^{-r})=1/r$. Their limsup contains every $c$ with infinitely many ones. Indeed, along its prefixes, $|R_i-c|_2=2^{-d_{i+1}}<1/R_i$. The remaining $c$ are countable. Haar measure is comparable to radius on all small balls in $\mathbb Z_2$, hence comparable to $\mathcal H^g$, and $h_1/g$ is monotone near zero. The cited theorem now gives $\mathcal H^{h_1}(E\cap U)=\mathcal H^{h_1}(U)=\infty$. For the last equality use $h_1(\rho)/\rho\to\infty$ and the positive Haar measure of $U$. For $0<s<1$ the gauge is larger near zero.
\end{proof}

\section{Davis's reciprocal-binomial conjectures}\label{dav:section}
Let $v=v_2$, let $s(n)$ count the ones in the binary expansion of a
nonnegative integer, and let $\ell(n)=\lfloor\log_2 n\rfloor$ for
$n>0$. Donald M. Davis formulated the following two estimates in his
2013 preprint, subsequently published in the \emph{Journal of
Combinatorics and Number Theory} (2014). His Theorem~2.2 proved an
outer-range estimate toward the second one. Both conjectures concern
the original rational numbers
\[
                    f(n)=\sum_{j=0}^n\binom nj^{-1}.
\]

\begin{theorem}[Davis's Conjectures 1.2 and 1.4]\label{dav:main}
For every integer $e\ge0$ and $0\le k<2^e$,
\begin{equation}\label{dav:first}
 v\bigl(f(2^e+k)-f(k)\bigr)\ge e-2s(k)-2.
\end{equation}
For every integer $e\ge1$ and $0\le k<2^{e-1}$,
\begin{equation}\label{dav:second}
 v\bigl(f(2^e+2k+1)-f(2k+1)\bigr)
                 \ge e-2\ell(k+3)+2v(k+1).
\end{equation}
Equality in \eqref{dav:first} holds exactly at the admissible values
$k=2^e-2$ and $k=2^e-4$. Equality in \eqref{dav:second} holds exactly at the admissible value $k=2^{e-1}-2$.
\end{theorem}

\subsection{The digit-sum estimate}
Combining \eqref{eq:tail} and \eqref{eq:pair} gives, for every $N=n+1$ and $t\ge1$,
\begin{equation}\label{dav:tail}
 f(n+t)-f(n)=-t\sum_{r\ge1}\frac{2^r r}{(N+r)(N+t+r)} .
\end{equation}
For \eqref{dav:first} take $n=k$ and $t=2^e$, and put
$A=v(N+r), B=v(N+t+r)$. The lowest $A$ digits of $k+r$
are ones. The lowest $B$ digits of $k+t+r$ are ones.
Binary digit sums are subadditive, and $k<t$ gives $s(k+t)=s(k)+1$.
Thus
\begin{equation}\label{dav:pop}
 A\le s(k)+s(r),\qquad B\le s(k)+1+s(r).
\end{equation}
Also $2s(r)\le r+1$: each set bit of weight at least $2$ pays for
two units, while the unit bit costs one additional unit. The $r$th
term in \eqref{dav:tail}, including $t$, has valuation
\[
 e+r+v(r)-A-B\ge e+r+v(r)-2s(k)-1-2s(r)\ge e-2s(k)-2.
\]
The non-Archimedean triangle inequality proves \eqref{dav:first}.

Equality for a summand forces equality in both digit-sum bounds and
$r+v(r)=2s(r)-1$. This last identity holds only for $r=1$ and $r=3$. Indeed $r$ must be odd and no set bit can have weight greater than $2$.
Equality in \eqref{dav:pop} then gives $B=A+1$, which forces $A=e$
because the two denominators differ by $2^e$. For $e\ge2$,
$0<N+r\le2^e+3<3\cdot2^e$ and hence $N+r=2^e$. For $e=0,1$
the possible $k$ and $r\in\{1,3\}$ can be inspected directly,
giving just $(e,k,r)=(1,0,1)$. Conversely, for each admissible
$k=2^e-2$ or $2^e-4$ the respective $r=1$ or $3$ attains the
bound ($s(k)=e-1$ or $e-2$ respectively). Every other summand
has larger valuation. This proves the equality classification.

\subsection{The odd-argument estimate}
Put $a=k+1$, $m=\ell(a+2)$, $h=v(a)$, so that
$N=2a\le 2^e=t$ and $m\ge h$. With $A=v(N+r)$ and
$B=v(N+t+r)$ it suffices to show, for every $r\ge1$,
\begin{equation}\label{dav:pole}
              A+B\le r+v(r)+2(m-h).
\end{equation}
For odd $r$ both valuations vanish. Write $r=2d$ and $z=v(d)$
otherwise. If $z<h$, then $A=B=1+z$ because $e\ge h+1>z+1$.
Here $2d+1+z\ge2+2z$, since $2d\ge z+1$.

For $z\ge h$ write $l=\ell(a+d)$. In every case
\[
 A+B\le 2l+3.
\]
Indeed, if $2a+2d<t$, then $A=B\le l+1$.
If $2a+2d\ge t$, its shift by $t$ has size at most $2(2a+2d)$,
so $A\le l+1$ and $B\le l+2$.
It is therefore enough to prove
\[
                       2(l-m)+h+2\le2d
\]
because $v(r)=1+z\ge1+h$. At $d=1$ we have $h=0$ and $l\le m$.
At $d=2$ we have $h\le1$ and $l=m$.
For $d\ge3$ set $b=\ell(d)$. The inequality
$a+d\le 2^b(a+2)$ gives $l-m\le b$.
For $d=3$ we have $h=0$ and $2b+2=4\le6$.
For $b\ge2$, use $h\le z\le b$ and
$3b+2\le2^{b+1}\le2d$. The first inequality follows by induction
starting at $b=2$. This proves \eqref{dav:pole}, and summing
\eqref{dav:tail} proves \eqref{dav:second}.

To attain equality in \eqref{dav:second}, one summand must have
the target valuation. For $r$ odd, the inequality is strict.
In the case $z<h$ we have $2d>z+1$, also giving strictness.
In the case $z\ge h$, the displayed bounds are strict for
$d=2,3$ and all $d\ge5$, since $3b+2<2d$ whenever
$d\ge5$. At $d=4$ use $a+4\le2(a+2)$ and $h\le2$ to
get $2(l-m)+h+2\le6<8$. This leaves $d=1$ and $h=0$.
Equality forces $l=m$, $A=l+1$ and $B=l+2$, since in the first
case $2a+2<t$ the bounds were strict. As $B>A$ and the
denominators differ by $2^e$, it follows that $A=e$.
The positive integer $2a+2\le2^e+2$ must therefore equal $2^e$.
Conversely, if $a=2^{e-1}-1\ge1$, then $h=0$, $m=e-1$,
$A=e$ and $B=e+1$ for $r=2$, giving equality for that unique
summand. This proves both equality statements observed experimentally
by Davis.

\begin{corollary}[Davis's $2$-definability criterion]\label{dav:cauchy}
Let $x_j$ be the least nonnegative residue of $x\in\mathbb Z_2$
modulo $2^j$. If $j-2s(x_j)\longrightarrow+\infty$, then the
sequence $f(x_j)$ converges in $\mathbb Q_2$.
\end{corollary}
\begin{proof}
A new bit at position $j$ changes the integer input by $2^j$.
At such a step \eqref{dav:first} bounds the valuation of the
increment below by $j-2s(x_j)-2$. All other increments vanish.
The Cauchy criterion in $\mathbb Q_2$ proves convergence.
\end{proof}
For instance, $x=\sum_{j\ge0}2^{3j}=-1/7$ satisfies the criterion.
Davis's Proposition~1.3 established this implication conditional on
Conjecture~1.2. The result now supplies the unconditional implication,
including all binary expansions with limiting frequency of ones
strictly less than $1/2$.


\section{Linear independence of the values at rational arguments}\label{tr:section}

Fix an embedding of
$\overline{\mathbb Q}$ into $\mathbb C_2$. For the rational arguments
used below, the dyadic tail gives
\begin{equation}\label{tr:tail}
             F(x)=-(x+1)\sum_{r\ge1}\frac{2^r}{x+1+r}.
\end{equation}
Indeed, if $x+1=a/q$ with $q$ odd and $a/q\notin\mathbb Z$, then
$v_2(x+1+r)=v_2(a+qr)\le\log_2|a+qr|=O(\log(2+r))$, so the
convergence criterion and the value formula apply.

\begin{theorem}\label{tr:main}
The set
\[
 \{1\}\ \cup \bigl\{F(x):x\in\mathbb Q\cap\mathbb Z_2,\ -1<x<0\bigr\}
\]
is linearly independent over $\overline{\mathbb Q}$. The inequalities
on $x$ refer to the ordinary real ordering of rational numbers.
Every $F(x)$ with $x\in(\mathbb Q\cap\mathbb Z_2)\setminus\mathbb Z$
is transcendental. For any nonempty finite set $S$ of such inputs,
\[
 \dim_{\overline{\mathbb Q}}
 \operatorname{span}_{\overline{\mathbb Q}}
       \bigl(\{1\}\cup\{F(x):x\in S\}\bigr)
       =1+\#\{x+\mathbb Z:x\in S\}.
\]
\end{theorem}

The multiplicative relations needed for the proof can be determined
before taking any $2$-adic logarithms.

\begin{lemma}\label{tr:units}
Let $q>1$ be odd, let $t$ be the positive real $q$th root of $2$,
and let $\zeta=\exp(2\pi i/q)$. Write $u_j=1-t\zeta^j$ for
$0\le j<q$. For integers $m_0,\ldots,m_{q-1}$, the product
$\prod_j u_j^{m_j}$ is a root of unity if and only if all $m_j$
are equal.
\end{lemma}
\begin{proof}
The identity $\prod_j u_j=1-t^q=-1$ gives one implication.
For the converse take $K=\mathbb Q(\zeta,t)$. The prime $2$ is
unramified in $\mathbb Q(\zeta)$ since $q$ is odd. At any prime
above $2$, the polynomial $X^q-2$ is therefore Eisenstein. Its $q$
roots $\zeta^it$ are conjugate over $\mathbb Q(\zeta)$, and each map
$t\mapsto\zeta^it$ extends to an automorphism of $K$ fixing
$\zeta$. All conjugates of a root of unity have absolute value one.
Setting $h_j=\log|1-t\zeta^j|$ with cyclic indices gives
\[
                \sum_{j=0}^{q-1}m_jh_{i+j}=0
                 \qquad(0\le i<q).
\]
The Fourier coefficient at zero of $h$ is $\log|1-t^q|=0$.
For $1\le b<q$, expand the ordinary real logarithm as
\[
 h_j=\log t-\sum_{n\ge1}\frac{t^{-n}}{n}
                  \cos(2\pi jn/q).
\]
The series is absolutely convergent. Orthogonality gives
\begin{equation}\label{tr:fourier}
 \sum_{j=0}^{q-1}\zeta^{-bj}h_j
 =-\frac q2\left(
 \sum_{\substack{n\ge1\\n\equiv b\ (q)}}\frac{t^{-n}}{n}
 +
 \sum_{\substack{n\ge1\\n\equiv -b\ (q)}}\frac{t^{-n}}{n}
 \right)<0 .
\end{equation}
Multiplying the preceding relations by $\zeta^{-bi}$ and summing
over $i$ yields
$(\sum_j m_j\zeta^{bj})(\sum_k h_k\zeta^{-bk})=0$.
Equation~\eqref{tr:fourier} forces every nonconstant Fourier
coefficient of $m$ to vanish. Fourier inversion makes $m$ constant.
\end{proof}

\begin{proof}[Proof of Theorem~\ref{tr:main}]
Now regard $K$ as embedded in $\mathbb C_2$, and retain $t,\zeta,u_j$
for the respective images. Since $|t|_2=2^{-1/q}<1$, all the
$u_j$ are principal units. Put $\lambda_j=\log u_j$, using the
$2$-adic logarithm in this proof. Their sum is $\log(-1)=0$.
Lemma~\ref{tr:units} shows that $u_0,\ldots,u_{q-2}$ are
multiplicatively independent. The inhomogeneous $p$-adic theorem on
linear forms in logarithms (Kaufman's estimate in the form of
\cite[Lemma~12 and Corollary~13]{MS}) gives
\begin{equation}\label{tr:baker}
                1,\lambda_0,\ldots,\lambda_{q-2}
      \quad\text{linearly independent over }\overline{\mathbb Q}.
\end{equation}
To check its smallness hypothesis, fix the number field containing
the units and the coefficients of a proposed relation and raise all
the units to the same sufficiently large power $2^h$. Their powers
approach $1$ because they belong to the principal units of a finite
extension of $\mathbb Q_2$. The degrees of the fields generated by the
powers and coefficients are bounded. The powers remain multiplicatively
independent and their logarithms are $2^h\lambda_j$. The cited corollary
makes every nontrivial algebraic linear combination of these logarithms transcendental, which gives~\eqref{tr:baker}.

Define convergent $2$-adic series for $0\le b<q$ by
\[
 P_b=\sum_{\substack{n\ge1\\ n\equiv b\ (q)}}\frac{t^n}{n}
     =-\frac1q\sum_{j=0}^{q-1}\zeta^{-bj}\lambda_j .
\]
The second equality follows by expanding $-\log(1-t\zeta^j)$
and using the root-of-unity filter. Since $P_0=0$ and the full
Fourier matrix is invertible, $P_1,\ldots,P_{q-1}$ span the same
space as $\lambda_0,\ldots,\lambda_{q-2}$. By~\eqref{tr:baker}
these $q-1$ elements together with $1$ are linearly independent.
For $1\le b<q$, separating the initial term in $P_b$ now gives
\begin{equation}\label{tr:shifted}
                F(-1+b/q)=1-bt^{-b}P_b .
\end{equation}
Indeed $t^{-b}P_b=\sum_{r\ge0}2^r/(qr+b)$, and the $r=0$ term is
$1/b$. This proves independence for the full set of fractions
with common denominator $q$, including the unreduced fractions.
Every finite subset of the displayed set in the theorem admits
such an odd common denominator.

Finally, reindexing the convergent series in~\eqref{tr:tail} gives
\[
             F(x+1)=1+\frac{x+2}{2(x+1)}F(x)
             \qquad(x\in(\mathbb Q\cap\mathbb Z_2)\setminus\mathbb Z).
\]
The rational multiplier is nonzero at each input in this formula.
A finite sequence of forward or backward integer shifts therefore
expresses each value as a rational affine transform, with nonzero
linear term, of the value at the unique representative in $(-1,0)$.
This proves both assertions about arbitrary rational nonintegral inputs.
\end{proof}

For example, the two original limits $F(-1/3)$ and $F(-2/3)$
are linearly independent together with $1$. The independence here
concerns algebraic linear combinations. Polynomial relations of
higher degree form a separate question.

For comparison, David, Hirata-Kohno and Kawashima \cite[Theorem~2.1]{DHK} give independence criteria for multiple polylogarithmic depths and rational shifts, under a place-dependent inequality. Here the weight is one and the argument is exactly $2$. Formula~\eqref{tr:fourier} calculates the full rank from the signs of the nonconstant Fourier coefficients.

\section{Antecedents and priority}\label{sec:priority}
The original targets have independent published formulations. Theorem~\ref{bbp:main} computes the kernel on the minimal generators of Borwein, Borwein and Galway's Definition~1. Theorem~\ref{bbg:log-exclusion} answers their Appendix D, Problem 8, and Corollary~\ref{bbp:log23} addresses the example singled out after their Theorem~6. Chamberland \cite{Ch} gives constructions of binary BBP logarithms for Gaussian-Mersenne primes. Borwein, Borwein and Galway's Gaussian valuations give the two low relations and permit finite searches. Their Appendix D also asks about general digit-extraction series, a broader class than the generators of Definition~1. The two-primitive-divisor result in \cite{AKNOT,Grech} supplies an input from finite group theory. Its use on the two Gaussian norms at the highest occupied level bounds the support of every putative relation. Adegoke, Lafont and Layeni give identities in a broader class of binary BBP series \cite{ALL}. Kulik discusses the dimension question for arrays of coefficients at fixed base and period \cite{Kulik}. Neither formulation gives the all-index kernel for the arctangents of Definition~1.

For Davis's sums, the earlier positive and negative cases include his Theorems~1.9, 1.12 and 4.10 \cite{Davis12}. His follow-up proves part (i) of Proposition~2.1 in its Theorem~2.2 and reports that both proposed bounds and their equality patterns were tested through $e=15$ \cite{Davis14}. Theorem~\ref{th:domain} handles both directions for every input, including the critical type radius $1/2$. Theorem~\ref{dav:main} proves the two finite inequalities with their original right-hand sides. The moving-index estimate \eqref{eq:badterm} and the uniqueness of each equality term are the noncancellation steps. Sections~\ref{sec:metric} and \ref{tr:section} concern consequences and further arithmetic questions.

One tempting extra claim has a definite predecessor. After Proposition~4.7 Davis conjectured $v(f(2^e-1))=3e-2$ for $e\ge4$. Lengyel's earlier Conjecture~3 concerns the equivalent truncated logarithm \cite{Lengyel}. Cohen's Theorem~1 gives its valuation as $2^e+2e-2$ \cite{Cohen}. Equation~\eqref{eq:finite} yields Davis's formula immediately. Cohen also completes the logarithmic tail and isolates a unique summand at these fixed inputs. We credit him with that calculation.

The literature check used the two author groups, the numbered questions, reciprocal-binomial and truncated-polylogarithm formulations, BBP zero relations and Gaussian primitive-divisor formulations. It also covered the type and rational-shift Lerch precedents \cite{Kedlaya,DHK}. These comparisons found no earlier full domain criterion, proof of both finite inequalities, full Aurifeuillian exclusion criterion, or all-index kernel theorem for the specified real generators. The search bounds the historical evidence for novelty. Independent bibliographic verification of priority remains open. All proof statements are independent of this historical judgment.

{\small\begin{thebibliography}{99}
\bibitem{ALL} K. Adegoke, J. Oliver Lafont and O. Layeni, A class of digit extraction BBP-type formulas in general binary bases, \emph{Notes on Number Theory and Discrete Mathematics} \textbf{17} (2011), no. 4, 18--32. \href{https://nntdm.net/volume-17-2011/number-4/18-32/}{Journal record}.
\bibitem{AKNOT} L. Alp\"oge, N. M. Katz, G. Navarro, E. A. O'Brien and P. H. Tiep, \emph{Local systems and Suzuki groups}, Theorem 3.4. \href{https://web.math.princeton.edu/~nmk/kt21_45bis.pdf}{Author's manuscript}.
\bibitem{BV} V. Beresnevich and S. Velani, A mass transference principle and the Duffin--Schaeffer conjecture for Hausdorff measures, \emph{Annals of Mathematics} \textbf{164} (2006), 971--992. \href{https://annals.math.princeton.edu/2006/164-3/p06}{Journal record}.
\bibitem{BBG} J. M. Borwein, D. Borwein and W. F. Galway, Finding and excluding $b$-ary Machin-type individual digit formulae, \emph{Canadian Journal of Mathematics} \textbf{56} (2004), 897--925. \href{https://doi.org/10.4153/CJM-2004-041-2}{doi:10.4153/CJM-2004-041-2}.
\bibitem{Ch} M. Chamberland, Binary BBP-formulae for logarithms and generalized Gaussian-Mersenne primes, \emph{Journal of Integer Sequences} \textbf{6} (2003), Article 03.3.7. \href{https://cs.uwaterloo.ca/journals/JIS/VOL6/Chamberland/chamberland60.pdf}{Publisher's PDF}.
\bibitem{Cohen} H. Cohen, On the 2-adic valuations of the truncated polylogarithm series, \emph{Fibonacci Quarterly} \textbf{37} (1999), 117--121. \href{https://www.fq.math.ca/Scanned/37-2/cohen.pdf}{Publisher's PDF}.
\bibitem{DHK} S. David, N. Hirata-Kohno and M. Kawashima, Linear independence criteria for generalized polylogarithms with distinct shifts, \emph{Acta Arithmetica} \textbf{206} (2022), 127--169. \href{https://arxiv.org/abs/2202.13931}{arXiv:2202.13931}.
\bibitem{Davis12} D. M. Davis, For which $p$-adic integers $x$ can $\sum_k\binom{x}{k}^{-1}$ be defined?, \emph{Journal of Combinatorics and Number Theory} \textbf{4} (2012), 87--104. \href{https://www.lehigh.edu/~dmd1/define3.pdf}{Author's PDF} and \href{https://arxiv.org/abs/1208.0250}{arXiv:1208.0250}.
\bibitem{Davis14} D. M. Davis, For which $2$-adic integers $x$ can $\sum_k\binom{x}{k}^{-1}$ be defined?, \emph{Journal of Combinatorics and Number Theory} \textbf{6} (2014), 215--218. \href{https://arxiv.org/abs/1301.2532}{arXiv:1301.2532} (2013).
\bibitem{Grech} M. A. Grechkoseeva, Primitive prime divisors of orders of Suzuki--Ree groups, \emph{Algebra and Logic} \textbf{62} (2023), 41--49. \href{https://arxiv.org/abs/2504.04729}{Corrected version, arXiv:2504.04729}, Theorem~2.
\bibitem{Kedlaya} K. S. Kedlaya, \emph{Exponents and $p$-adic Liouville numbers}, course 18.787, Chapter 12, Definition~12.1.1 and Propositions~12.1.4--5. \href{https://kskedlaya.org/18.787/exponents.pdf}{Author's notes}.
\bibitem{Kulik} D. Kulik, \emph{The Reader and the Frame}, 2026, Section~10. \href{https://doi.org/10.5281/zenodo.22007472}{Zenodo record}.
\bibitem{Lengyel} T. Lengyel, Characterizing the 2-adic order of the logarithm, \emph{Fibonacci Quarterly} \textbf{32} (1994), 397--401. \href{https://www.fq.math.ca/Scanned/32-5/lengyel.pdf}{Publisher's PDF}.
\bibitem{MS} M. Ram Murty and N. Saradha, Transcendental values of the $p$-adic digamma function, \emph{Acta Arithmetica} \textbf{133} (2008), 349--362. \href{https://doi.org/10.4064/aa133-4-4}{doi:10.4064/aa133-4-4}.
\bibitem{Rockett} A. M. Rockett, Sums of the inverses of binomial coefficients, \emph{Fibonacci Quarterly} \textbf{19} (1981), 433--437. \href{https://www.fq.math.ca/Scanned/19-5/rockett.pdf}{Publisher's PDF}.
\bibitem{Corpus} QEEDDD, \emph{Verifiable Mathematical Reasoning Corpus}, items 022, 139, 169, 176 and 180. \href{https://github.com/QEEDDD/Verifiable-Mathematical-Reasoning-Corpus/tree/5887cdea4898f0a3aa5bf7a8e0829697bd2b0f4b}{Repository snapshot 5887cdea}.
\end{thebibliography}}
\end{document}
